\begin{afterword}
\summaryhead

The post states Newcomb's problem. A predictor with a perfect record leaves \$1,000 in a transparent box and puts \$1,000,000 in an opaque box only if it predicts that you will take the opaque box alone. The two-boxer argues that the boxes are already filled, so taking both gains \$1,000 either way. The one-boxer asks, ``why ain'cha rich?''

The author grants that causal decision theory, which recommends two boxes, is the dominant view, and that even its defenders agree one should precommit to one box if one can. The author's own theory is withheld, and the post gives motivations instead. Rational agents should win. Omega rewards the predicted choice, not any way of reasoning, so it does not reward irrationality. Joyce's causal decision theorist admits wishing to be the one-boxing type; the author replies that one should not envy another's mere choice, but make it. If one would wish that one-boxing were reasonable, one should revise ``reasonable.'' The post concedes that this is not a knockdown criticism.

\responsehead

The post is open about what it is. It states the two-boxer's case fairly, grants that the opposing view is the dominant one, quotes a standard causal reply at length in its own words, and ends by conceding that it offers motivations, not a refutation. The surveys confirm its picture of the field. In 2009 and 2020, more philosophers chose two boxes than one, and in 2020 decision theorists were more strongly for two boxes than other philosophers. The one-boxing side also had an established theory before the post: evidential decision theory, which picks the act that is the best evidence of a good outcome.

The post's thesis, that rational agents should win, is common ground. Causal decision theorists, who pick the act that causes the best outcome, also want the money, and argue that two boxes get \$1,000 more of it whatever box B already holds. The dispute is over which comparison defines winning, and the post answers it from the one-boxer's side; it presents its arguments as ``motivations,'' not as the theory it claims to have. Lewis, who put ``Why ain'cha rich?'' in a title, concluded in 1981 that on whether the one-boxers' riches prove anything, ``it's a standoff.'' The notes give the causal reply to each of the post's motivations.

One objection the post answers is given without a source: that one-boxers must believe their choice affects the box. Lewis's charge is different: the evidential rule seeks good news about matters the choice cannot affect. That charge needs no belief that the choice affects the box.

In short: a candid statement of the one-boxing side that reports the opposing view fairly and offers, as it admits, motivations rather than a refutation of the causal reply; the dispute remains, in Lewis's word, a standoff.
\end{afterword}
